\chapter{Décrire un projet : sources et sorties}

Un projet se décrit par ce qu'il \emph{est}, ce qu'il \emph{contient} et ce qu'il
\emph{produit}. Ce chapitre couvre les trois, avec les valeurs par défaut --- car
la moitié du travail consiste à ne rien écrire.

\section{Ce qu'un projet EST : son genre}

\begin{center}
\begin{tabular}{@{}lll@{}}
\toprule
\textbf{Dans le \texttt{.jenga}} & \textbf{En ligne de commande} & \textbf{Produit} \\
\midrule
\texttt{consoleapp()} & \texttt{-{}-kind console} & un exécutable avec console \\
\texttt{windowedapp()} & \texttt{-{}-kind windowed} & un exécutable sans console \\
\texttt{staticlib()} & \texttt{-{}-kind static} & une bibliothèque statique \\
\texttt{sharedlib()} & \texttt{-{}-kind shared} & une bibliothèque dynamique \\
\texttt{testsuite()} & \texttt{-{}-kind test} & une suite de tests \\
\bottomrule
\end{tabular}
\end{center}

\begin{nkretenir}
\textbf{\texttt{windowedapp()} n'est pas cosmétique sous Windows.} Il change le
sous-système de l'exécutable : sans lui, une fenêtre s'ouvre avec une console
noire derrière elle.

Et le point d'entrée attendu change avec le sous-système --- c'est ce qui produit
le fameux \texttt{undefined reference to WinMain} quand la déclaration et le code
ne s'accordent pas.
\end{nkretenir}

\begin{nkretenir}
\textbf{Le genre se déclare souvent DEUX fois}, et ce n'est pas une erreur :
une fois au niveau du projet, une fois dans chaque filtre de plateforme.
Certaines plateformes en changent --- le Web est déclaré \texttt{consoleapp()}
là où le bureau est \texttt{windowedapp()}, parce que la notion de fenêtre
native n'y a pas le même sens.
\end{nkretenir}

\section{Le langage et son dialecte}

\begin{nkcode}{}
language("C++")
cppdialect("C++17")
\end{nkcode}

\begin{center}
\begin{tabular}{@{}ll@{}}
\toprule
\textbf{Fonction} & \textbf{Équivalent} \\
\midrule
\texttt{cppdialect(\dots)} & \texttt{cppversion(\dots)} \\
\texttt{cdialect(\dots)} & \texttt{cversion(\dots)} \\
\bottomrule
\end{tabular}
\end{center}

\begin{nkretenir}
Les deux noms font la même chose. \textbf{Choisissez-en un et tenez-vous-y} :
deux orthographes pour un même réglage rendent une recherche dans le dépôt
incomplète --- on cherche \texttt{cppdialect} et l'on rate les fichiers qui
disent \texttt{cppversion}.
\end{nkretenir}

\section{Où le projet vit : \texttt{location}}

\begin{nkcode}{}
location(".")          # le projet vit dans le dossier de son .jenga
location("Bonjour")    # ... dans un sous-dossier
\end{nkcode}

\begin{nkretenir}
\textbf{Tous les chemins du projet sont relatifs à \texttt{location}}, et
\texttt{location} est lui-même relatif au dossier du \texttt{.jenga} --- puisque
le chargeur s'y place (chapitre 4).

Un descripteur d'application inclus depuis la racine écrit donc
\texttt{location(".")} et parle de son propre dossier, pas de la racine.
\end{nkretenir}

\section{Quels fichiers lui appartiennent}

\begin{nkcode}{}
files(["src/**.cpp"])                        # recursif
files(["src/*.cpp", "src/Ui/*.cpp"])         # non recursif, deux dossiers
excludefiles(["src/vieux/**.cpp"])
\end{nkcode}

\begin{center}
\begin{tabular}{@{}ll@{}}
\toprule
\textbf{Motif} & \textbf{Portée} \\
\midrule
\texttt{*.cpp} & le dossier courant seulement \\
\texttt{**.cpp} & le dossier et tous ses sous-dossiers \\
\bottomrule
\end{tabular}
\end{center}

\begin{nkpiege}
\textbf{Un motif ne dit pas ce qu'il ne trouve pas.} \texttt{files(["src/**.cpp"])}
sur un dossier vide donne un projet à zéro source --- et l'échec sort à l'édition
de liens, sur un symbole manquant, sans jamais mentionner que rien n'a été
compilé.

La sortie de construction porte pourtant la réponse :
\texttt{Found 1 source file(s)}. \textbf{Lisez ce nombre} ; c'est la ligne qui
vous dit que votre motif a mordu.
\end{nkpiege}

\begin{nkretenir}
\texttt{excludemainfiles()} existe pour un cas précis : retirer les fichiers qui
portent un \texttt{main()} quand on veut réutiliser les mêmes sources dans une
bibliothèque ou une suite de tests. Deux \texttt{main} dans une même édition de
liens ne pardonnent pas.
\end{nkretenir}

\section{Ce qu'il produit, et où}

\begin{center}
\begin{tabular}{@{}ll@{}}
\toprule
\textbf{Fonction} & \textbf{Ce qu'elle fixe} \\
\midrule
\texttt{objdir(chemin)} & où vont les fichiers objets \\
\texttt{targetdir(chemin)} & où va l'artefact final \\
\texttt{targetname(nom)} & le nom de l'artefact, si différent du projet \\
\bottomrule
\end{tabular}
\end{center}

\begin{nkcode}{La forme employee partout dans Nkentseu}
objdir("%{wks.location}/Build/Obj/%{cfg.buildcfg}-%{cfg.system}/%{prj.name}")
targetdir("%{wks.location}/Build/Bin/%{cfg.buildcfg}-%{cfg.system}/%{prj.name}")
\end{nkcode}

\begin{nkretenir}
\textbf{Les \texttt{\%\{\dots\}} sont des variables résolues par Jenga}, pas du
Python :

\begin{center}
\begin{tabular}{@{}ll@{}}
\toprule
\texttt{\%\{wks.location\}} & le dossier du workspace \\
\texttt{\%\{cfg.buildcfg\}} & \texttt{Debug} ou \texttt{Release} \\
\texttt{\%\{cfg.system\}} & \texttt{Windows}, \texttt{Linux}\dots \\
\texttt{\%\{prj.name\}} & le nom du projet \\
\bottomrule
\end{tabular}
\end{center}

\textbf{Sans la configuration et le système dans le chemin}, un build Debug
écraserait le Release et un build Linux écraserait le Windows. Ce n'est pas du
rangement : c'est ce qui rend deux constructions simultanées possibles.
\end{nkretenir}

\begin{nkpiege}
\textbf{Certaines plateformes ajoutent un suffixe.} Sous Linux, les descripteurs
de Nkentseu déclarent un dossier par backend graphique :

\begin{nkcode}{Applications/NkDames/NkDames.jenga}
objdir(".../Build/Obj/%{cfg.buildcfg}-%{cfg.system}-xlib/%{prj.name}")
\end{nkcode}

D'où \texttt{Release-Linux-xlib} et non \texttt{Release-Linux}. Un script
d'empaquetage qui code le second en dur échoue sur une construction
\textbf{réussie} --- et c'est arrivé.
\end{nkpiege}

\section{Ce qu'on ne déclare pas}

\begin{nkretenir}
\textbf{La moitié d'un bon descripteur est ce qu'on n'y écrit pas.} Le
compilateur, le niveau d'avertissement, l'extension de l'artefact, la façon
d'appeler l'éditeur de liens : Jenga les déduit du genre, du langage et de la
plateforme.

\textbf{N'écrivez un réglage que lorsque le défaut ne convient pas} --- et quand
vous le faites, écrivez \emph{pourquoi} à côté. Un drapeau sans justification
survit dix ans après que sa raison a disparu.
\end{nkretenir}

\section{Les métadonnées de l'application}

\begin{nkcode}{}
apppublisher("Rihen Universe")
appversion("0.1.0")
licensefile("../../LICENSE")
appicon("Resources/icone.png")
\end{nkcode}

\begin{nkretenir}
Elles ne servent pas à la compilation mais à l'\emph{empaquetage} : elles
remplissent le manifeste Android, les propriétés de l'exécutable Windows,
l'installeur.

\textbf{\texttt{appicon} pointe un fichier qui doit exister.} Déclarer une icône
absente fait échouer la génération de l'APK --- alors que la construction bureau,
elle, ne dit rien. Le descripteur de NkDames porte d'ailleurs le commentaire, et
la ligne est mise en attente jusqu'à ce que l'icône soit dessinée.
\end{nkretenir}

\begin{nkbilan}
Un projet déclare son genre, son langage, son dossier, ses fichiers et ses
sorties --- et rien de plus si les défauts conviennent. Le genre change le
sous-système et donc le point d'entrée attendu ; les motifs \texttt{*} et
\texttt{**} ne fouillent pas de la même façon, et la ligne
\texttt{Found N source file(s)} est ce qui vous dit si le vôtre a mordu. Les
chemins de sortie portent la configuration et le système parce que sans eux deux
constructions s'écrasent --- et certaines plateformes y ajoutent un suffixe qu'un
script ne doit pas coder en dur. Enfin les métadonnées ne servent qu'à
l'empaquetage, mais une icône déclarée doit exister.
\end{nkbilan}

\begin{nkquestions}
\begin{enumerate}
\item Quelle différence entre \texttt{consoleapp()} et \texttt{windowedapp()},
      au-delà de la console ?
\item Que fouille \texttt{src/*.cpp} que \texttt{src/**.cpp} fouille aussi, et
      inversement ?
\item Quelle ligne de la sortie de construction dit qu'un motif a trouvé quelque
      chose ?
\item Pourquoi la configuration et le système figurent-ils dans le chemin de
      sortie ?
\item À quoi servent \texttt{apppublisher} et \texttt{appversion} ?
\end{enumerate}
\end{nkquestions}

\begin{nkpratique}
Reprenez le projet du chapitre 3 et faites-lui produire, tour à tour :

\begin{enumerate}
\item une bibliothèque statique au lieu d'un exécutable ;
\item un exécutable dont le \emph{nom} diffère du projet ;
\item un exécutable dont les objets vont dans un dossier à vous.
\end{enumerate}

Après chaque changement, \textbf{trouvez l'artefact sur le disque} avant de
relancer \texttt{jenga run}. C'est la seule façon de vérifier que le réglage a
été honoré et pas seulement accepté.
\end{nkpratique}

\begin{nkdemo}
Prouvez qu'un motif qui ne trouve rien ne le dit pas.

Créez un projet dont \texttt{files()} pointe un dossier \textbf{vide}. Construisez.

\textbf{Ce que la démonstration doit établir}, et c'est en deux temps :
\begin{enumerate}
\item relevez ce qu'affiche la ligne \texttt{Found \dots\ source file(s)} ;
\item relevez le message d'échec final, et constatez qu'il parle d'un
      \emph{symbole}, pas d'un motif.
\end{enumerate}

Puis ajoutez un fichier et refaites la mesure. C'est la paire qui prouve que la
première ligne était le seul indice utile.
\end{nkdemo}

\begin{nklogique}
Un collègue a écrit \texttt{targetdir("Build/Bin")} --- sans configuration ni
système.

\begin{enumerate}
\item Décrivez précisément ce qui se passe quand il construit en Debug, puis en
      Release, puis relance son programme.
\item Le symptôme qu'il rapportera n'est pas « mes builds s'écrasent ». Que
      dira-t-il, à votre avis, et pourquoi cette formulation envoie-t-elle
      chercher au mauvais endroit ?
\item Une troisième construction, pour une \emph{autre plateforme}, aggrave le
      cas d'une façon nouvelle. Laquelle ?
\end{enumerate}
\end{nklogique}
